\chapter{Practical applications: wiki and CMS}

The next two examples use the same foundations but demonstrate different product needs. The minimal wiki combines stable slugs, Markdown content, and authenticated editing; a more mature wiki can later add optimistic locking and an audit trail. The CMS example adds a multi-file project structure, editor authorization, and public cache, then lists the additional layers a production CMS usually needs.

The minimal wiki is also checked as a compilation fixture in \code{examples/wiki/main.rw}; the CMS is the multi-file \code{examples/news-site/} project. These examples are therefore not merely book pseudocode.

\section{A simple wiki}

A minimal wiki domain consists of pages: stable identifier, slug, title, and Markdown body. The following example focuses on connecting the language elements.

\subsection*{Model}

\begin{lstlisting}[language=RWLang]
model WikiPage {
    id: Int
    slug: Slug
    title: String
    body: String
}
\end{lstlisting}

\subsection*{Query}

\begin{lstlisting}[language=RWLang]
query fn wikiBySlug(db: Db, slug: Slug)
    -> Result<WikiPage, DbError> sql {
    SELECT id, slug, title, body
    FROM wiki_pages
    WHERE slug = :slug
}
\end{lstlisting}

\subsection*{Public page}

\begin{lstlisting}[language=RWLang]
page fn wikiShow(ctx: PageContext, db: Db, slug: Slug)
    -> Result<Html, PageError> {
    let page = wikiBySlug(db, slug)?;
    return Ok(html {
        @layout(Main, page.title) {
            <article>
                <h1>{{ page.title }}</h1>
                @markdown(page.body)
            </article>
        }
    });
}

route wikiShow GET "/wiki/:slug<Slug>" => wikiShow;
\end{lstlisting}

\subsection*{Editor form}

\begin{lstlisting}[language=RWLang]
form WikiForm {
    title<String>
    body<String>
    validate title length 2 200 body length 1 300000
}
\end{lstlisting}

The mutating route should require authentication, for example \code{auth role Editor}. The write query should receive a \code{Transaction} capability. Do not blindly regenerate and overwrite the slug from the title on every edit if preserving stable old URLs matters.

\subsection*{Canonical slug and aliases}

A more mature wiki can maintain old names in a separate slug-alias/history table. The runtime can produce a 301 redirect from a canonical route, so the application does not need to build arbitrary redirect URLs.

\subsection*{What should a real wiki add?}

\begin{itemize}
  \item optimistic locking for concurrent edits;
  \item a business audit trail for changes;
  \item object-level authorization if not every page has identical access rules;
  \item backup/restore policy for the database and any AppFs state.
\end{itemize}

\section{CMS and news site}

The repository's \code{examples/news-site/} project already demonstrates a small CMS core: model, typed SQL, layout, Markdown, editor action, and public cache.

\subsection*{Project structure}

\begin{lstlisting}
main.rw
models.rw
queries.rw
components.rw
pages.rw
actions.rw
\end{lstlisting}

\subsection*{Article model}

\begin{lstlisting}[language=RWLang]
model Article {
    id: Int
    slug: Slug
    title: String
    body: String
}
\end{lstlisting}

Editor input gets a separate named form schema:

\begin{lstlisting}[language=RWLang]
form ArticleForm {
    title<String>
    body<String>
    validate title length 2 200 body length 1 300000
}
\end{lstlisting}

\subsection*{Public article page}

\begin{lstlisting}[language=RWLang]
page fn article(ctx: PageContext, db: Db, slug: Slug)
    -> Result<Html, PageError> {
    let article = articleBySlug(db, slug)?;
    return Ok(html {
        @layout(Main, article.title) {
            <article>
                <h1>{{ article.title }}</h1>
                @markdown(article.body)
            </article>
        }
    });
}

route article GET "/article/:slug<Slug>"
    cache public ttl 60
    => article;
\end{lstlisting}

Public cache is appropriate only for a page that does not depend on session or user state.

\subsection*{Editor-side creation}

\begin{lstlisting}[language=RWLang]
action fn create(ctx: ActionContext, db: Db,
                 title: String, body: String)
    -> Result<Redirect, PageError> {
    let articleSlug = slug(title);
    transaction db {
        insertArticle(tx, articleSlug, title, body)?;
    }
    return Ok(redirect("/admin/articles"));
}

route articleCreate POST "/admin/articles"
    form ArticleForm
    auth role Editor
    invalidate cache article
    => create;
\end{lstlisting}

The invalidation above changes the route-level generation: it does not delete just one cache key for the new slug, but makes the previous generation of the \code{article} route unreachable. This is correct and predictable in a simple example; in a high-traffic CMS you should deliberately decide which public routes a given mutation really needs to invalidate.

\subsection*{From CMS to production system}

The following additions form the backbone of a practical production CMS:

\begin{itemize}
  \item separate validation for publication and domain invariants;
  \item draft/published state as an enum;
  \item canonical slug history;
  \item optimistic locking during editing;
  \item business audit trail;
  \item media library and AppFs for images;
  \item role-based editor/publisher policy;
  \item rate limits and resource profiles for expensive routes;
  \item backups, restore drills, and controlled deployment.
\end{itemize}

In a CMS, original Markdown should remain the source of truth; do not store raw rendered HTML produced by the application.
